\documentclass[12pt]{article}
\usepackage[utf8]{inputenc}
\usepackage[T1]{fontenc}
\usepackage{hyperref}
\title{Waka Language Resources for K-12 Curriculum Development}
\author{Compiled by Claude Sonnet 5 (claude-sonnet-5) for Andr\'as Kornai}
\date{}
\begin{document}
\maketitle
\section{Introduction}
Waka (ISO 639-3: wav; Glottocode waka1275) is a small Nigerian language of the Adamawa branch within Atlantic-Congo, with roughly 5,000 speakers cited. Sources disagree on where it is spoken (Taraba or Adamawa State) and Joshua Project gives larger population figures. Very little educational material was found.
\section{Primary Educational Resources}
No primers, textbooks or school programmes for Waka were found in search results.
\section{Linguistic and Reference Materials}
\begin{itemize}
\item Wikipedia overview of the language: \url{https://en.wikipedia.org/wiki/Waka_language}
\item Glottolog entry (classification): \url{https://glottolog.org/resource/languoid/id/waka1275.bigmap.html}
\item Numerals dataset page (incomplete beyond ten): \url{https://lingweb.eva.mpg.de/channumerals/Waka.htm}
\end{itemize}
\section{Religious and Community Resources}
Joshua Project language page (population figures only): \url{https://joshuaproject.net/languages/wav}
\section{Digital and Audio Resources}
None found.
\section{Curriculum Development Context}
Materials for a K-12 curriculum would have to be created from scratch. Note that the Kambari languages (Kainji, Kebbi and Niger States) are a different group and must not be confused with Waka.
\end{document}